(1) 介质覆盖的横截面积为 \(wx\)，其余真空面积 \(w(L-x)\)，两块区域在同一个电压下并联。故 \(\boxed{C(x)=\frac{\varepsilon_0w}{d}[L+(\kappa-1)x]}\)、\(C'(x)=\varepsilon_0w(\kappa-1)/d>0\)。无需错误地将沿板长的区域串联。

(2) 固定 \(V\) 时介质、场与电源共同的适当势为 \(\mathcal F=-C(x)V^2/2\)，故力 \(F_x=-\mathcal F'=V^2C'(x)/2>0\)，趋向吸入。电场能 \(U_E=C V^2/2\) 增加 \(\Delta U_E=V^2\Delta C/2\)；电源向系统提供 \(W_{\rm src}=V\Delta Q=V^2\Delta C\)。剩余一半是介质获得的机械功。

(3) 断开电源后固定自由电荷 \(Q\)，静电能为 \(U_E=Q^2/[2C(x)]\)。吸入力 \(\boxed{F_x=-dU_E/dx=Q^2C'(x)/[2C(x)^2]}>0\)，其机械功来自电场能下降，没有外部电源供功。若 \(Q=CV\) 在某一瞬间两种装置具有相同瞬时场和力，但随 \(x\) 改变的力学曲线不同。

(4) 当 \(\kappa\to1\) 时 \(C'\to0\)，介质不产生吸入力。理想区间 \(0<x<L\) 内力近似常数（固定电压），端点附近边缘场与实际几何使其变化；即使本题忽略大部分边缘场，正是插入前沿附近的非均匀电场产生净水平力，不能在无限完全均匀场里把其当成体积力。
